% 2010 年浙江卷文科第 8 题：三视图（正视图、侧视图、俯视图）。
% 来源 PDF：普通高考/2010/2010浙江文.pdf，SHA-256 a73e22ea405697775b2d5cd82532fcffad0f506b849d9f199265b4297b8a754b。
% 原页：第 1 页（A3 三栏）第 2 栏上部，书页码 641。
% 重绘依据：正视图与侧视图均为下底 8、上底 4、高 2 的等腰梯形，其上叠 4×2 的矩形；
%   俯视图为外方 8×8、内方 4×4、四角以对角线相连。尺寸 2、4、2 与 2、2 照原页标注。
\begin{tikzpicture}[scale=0.45, semithick, every node/.style={inner sep=1.5pt, font=\small},
  dim/.style={<->, >=stealth, thin}]

  % ---------- 正视图 ----------
  \draw (-4,0) -- (4,0) -- (2,2) -- (-2,2) -- cycle;
  \draw (-2,2) rectangle (2,4);
  % 顶部尺寸 2,4,2
  \draw[thin] (-4,4.25) -- (-4,4.85);
  \draw[thin] (-2,4.25) -- (-2,4.85);
  \draw[thin] (2,4.25) -- (2,4.85);
  \draw[thin] (4,4.25) -- (4,4.85);
  \draw[dim] (-4,4.55) -- (-2,4.55);
  \draw[dim] (-2,4.55) -- (2,4.55);
  \draw[dim] (2,4.55) -- (4,4.55);
  \node[above] at (-3,4.55) {$2$};
  \node[above] at (0,4.55) {$4$};
  \node[above] at (3,4.55) {$2$};
  % 右侧尺寸 2,2
  \draw[thin] (4.15,4) -- (4.75,4);
  \draw[thin] (4.15,2) -- (4.75,2);
  \draw[thin] (4.15,0) -- (4.75,0);
  \draw[dim] (4.45,2) -- (4.45,4);
  \draw[dim] (4.45,0) -- (4.45,2);
  \node[right] at (4.45,3) {$2$};
  \node[right] at (4.45,1) {$2$};
  \node[below] at (0,-0.9) {正视图};

  % ---------- 侧视图 ----------
  \begin{scope}[shift={(11,0)}]
    \draw (-4,0) -- (4,0) -- (2,2) -- (-2,2) -- cycle;
    \draw (-2,2) rectangle (2,4);
    \node[below] at (0,-0.9) {侧视图};
  \end{scope}

  % ---------- 俯视图 ----------
  \begin{scope}[shift={(0,-3.4)}]
    \draw (-4,0) rectangle (4,-8);
    \draw (-2,-2) rectangle (2,-6);
    \draw (-4,0) -- (-2,-2);
    \draw (4,0) -- (2,-2);
    \draw (4,-8) -- (2,-6);
    \draw (-4,-8) -- (-2,-6);
    % 右侧尺寸 2,4,2
    \draw[thin] (4.15,0) -- (4.75,0);
    \draw[thin] (4.15,-2) -- (4.75,-2);
    \draw[thin] (4.15,-6) -- (4.75,-6);
    \draw[thin] (4.15,-8) -- (4.75,-8);
    \draw[dim] (4.45,-2) -- (4.45,0);
    \draw[dim] (4.45,-6) -- (4.45,-2);
    \draw[dim] (4.45,-8) -- (4.45,-6);
    \node[right] at (4.45,-1) {$2$};
    \node[right] at (4.45,-4) {$4$};
    \node[right] at (4.45,-7) {$2$};
    \node[below] at (0,-8.9) {俯视图};
  \end{scope}
\end{tikzpicture}
